\section*{Supplementary Material}
\addcontentsline{toc}{section}{Supplementary Material}
\renewcommand{\thesection}{S\arabic{section}}
\setcounter{section}{0}

\section{Full Layer-wise Results (Qwen2.5-1.5B)}
\label{sup:layerwise}

\begin{table}[h]
\centering
\small
\caption{Per-layer SST-2 AUC and EmoBank V $|r|$ for Qwen2.5-1.5B (all 28
layers). The V-shape pattern: strong at L1, near-zero from L2 to L26, abrupt
jump at L27.}
\label{tab:sup_layerwise}
\begin{tabular}{lcc}
\toprule
Layer & SST-2 AUC & EmoBank V $|r|$ \\
\midrule
L1  & 0.752 & 0.370 \\
L2  & 0.524 & 0.005 \\
L3  & 0.524 & 0.005 \\
L4  & 0.524 & 0.005 \\
L8  & 0.524 & 0.005 \\
L12 & 0.525 & 0.005 \\
L14 & 0.525 & 0.005 \\
L16 & 0.525 & 0.005 \\
L18 & 0.525 & 0.005 \\
L20 & 0.525 & 0.005 \\
L22 & 0.526 & 0.004 \\
L24 & 0.527 & 0.004 \\
L25 & 0.531 & 0.003 \\
L26 & 0.532 & 0.002 \\
\textbf{L27} & \textbf{0.836} & \textbf{0.517} \\
\textbf{L28} & \textbf{0.837} & \textbf{0.496} \\
\bottomrule
\end{tabular}
\end{table}

\section{Latent Space Analysis}
\label{sup:latent}

We computed a set of geometric diagnostics per layer to quantify how the
residual stream changes across the network:

\begin{table}[h]
\centering
\small
\caption{Latent-space metrics per layer for Qwen2.5-1.5B. Isotropy is the
von-Neumann entropy of eigenvalue distribution (0 = fully anisotropic, 1 =
perfectly isotropic). Effective rank of the 9 centroids. Silhouette score
of the 9-class clustering. V-axis cosine similarity with the L28 V-axis.}
\label{tab:sup_latent}
\begin{tabular}{lccccc}
\toprule
Layer & Isotropy & Eff.\ rank & PC1 var ratio & Silhouette & cos$(v_L, v_{28})$ \\
\midrule
L1 & 0.010 & 2.99 & 0.972 & $-$0.178 & 0.291 \\
L14 & 0.033 & 3.94 & 0.937 & $-$0.145 & \textbf{0.067} \\
L27 & 0.576 & 6.92 & 0.515 & $-$0.008 & 0.953 \\
L28 & 0.533 & 6.80 & 0.546 & $-$0.014 & 1.000 \\
\bottomrule
\end{tabular}
\end{table}

Isotropy rises from 0.01 to 0.58 across layers---the residual stream becomes
increasingly high-dimensional / balanced across dimensions. The effective
rank of the 9 centroids roughly doubles (3.0 $\to$ 6.9). The silhouette
score is always negative, meaning individual story embeddings don't cluster
tightly by class in any layer---the 9-class structure is a
\emph{centroid-level} property, not a sample-level property.

\section{Multilingual Self-LOOCV}
\label{sup:multilingual}

Since all Hugging Face multilingual sentiment datasets (tyqiangz/multilingual-sentiments,
cardiffnlp/tweet\_sentiment\_multilingual) rely on deprecated script-based
loaders, we cannot test cross-language generalization on external
benchmarks directly. Instead, we perform \emph{self-LOOCV}: generate 20
stories per class in each language, use 12 for centroids and 8 for test,
and classify test stories by nearest centroid.

\begin{table}[h]
\centering
\caption{Multilingual self-LOOCV accuracy on Qwen2.5-1.5B. 9-class
nearest-centroid and binary (pos vs neg emotions) AUC.}
\label{tab:sup_multilingual}
\begin{tabular}{lccc}
\toprule
Language & 9-class acc & Binary AUC & Binary acc \\
\midrule
English & 0.764 & 0.535 & 0.563 \\
Spanish & 0.681 & 0.899 & 0.781 \\
French  & 0.597 & 0.658 & 0.594 \\
German  & 0.486 & 0.870 & 0.781 \\
Chinese & \textbf{0.875} & \textbf{0.967} & \textbf{0.875} \\
\bottomrule
\end{tabular}
\end{table}

The method works across all five tested languages. Binary valence AUC
is strongest in Chinese (0.97) and Spanish (0.90), moderate in German
(0.87) and French (0.66), and surprisingly weak on the small
English held-out self-LOOCV set (0.54); the same Qwen2.5-1.5B model
achieves AUC 0.868 on the independent SST-2 test set (\cref{tab:sst2}),
so the low EN-self-LOOCV AUC reflects sampling variance on only 72 held-out
stories, not a genuine weakness. Chinese also has the highest 9-class
nearest-centroid accuracy (0.875).

\section{V-shape is causal-LM specific: BERT and RoBERTa baselines}
\label{sup:bidirectional}

The V-shape pattern we identify (dilution in middle layers, crystallization at
the final 1--2 blocks) is specific to \emph{causal} (autoregressive) language
models. We verify this by running the same nine-story-centroid PCA method on
two bidirectional encoder-only models: BERT-base and RoBERTa-base.

\begin{table}[h]
\centering
\caption{Per-layer zero-shot SST-2 AUC for bidirectional encoders. Unlike
Qwen / Pythia / Bloom (\cref{fig:crossmodel}), the peak AUC occurs in
\emph{middle} layers and the final layer is \emph{not} the best.}
\label{tab:sup_bidirectional}
\begin{tabular}{lcccccccc}
\toprule
Model & L1 & L3 & L5 & L7 & L9 & L11 & L12 & Peak layer \\
\midrule
BERT-base      & 0.730 & 0.749 & 0.755 & 0.789 & \textbf{0.845} & 0.805 & 0.814 & L9  \\
RoBERTa-base   & 0.723 & 0.761 & 0.834 & \textbf{0.856} & 0.831 & 0.835 & 0.771 & L7  \\
\midrule
Qwen2.5-1.5B (causal, final-layer peak) & 0.752 & \ldots & 0.525 & \ldots & 0.525 & \ldots & 0.532 & \textbf{L28} \\
\bottomrule
\end{tabular}
\end{table}

BERT peaks at layer 9 of 12 (AUC 0.845) and RoBERTa at layer 7 of 12 (AUC
0.856). In both cases, the final layer has lower AUC than several middle
layers. Compare this to Qwen2.5-1.5B where layers 2--26 are all near chance
and the final 1--2 layers jump by 0.3 AUC.

\paragraph{Interpretation.}
Bidirectional encoders do not have to produce a next-token distribution at
the final layer, so there is no pressure to ``re-linearize'' semantic
features at the end. The V-axis lives in middle layers where contextual
mixing has surfaced valence-relevant features but late-layer compression
has not yet begun. Causal LMs, by contrast, \emph{must} surface decodable
next-token information in the final residual stream, which incidentally
re-linearizes valence. This is a testable prediction in the main text
(\cref{sec:discussion}) that we now verify empirically.

\section{Mechanism: Where is the V-axis built?}
\label{sup:mechanism}

The granular layer sweep (\cref{tab:sup_layerwise}) shows the V-axis
``crystallizes'' between paper layer L26 (SST-2 AUC 0.532) and L27
(AUC 0.836). To localize this within the transformer block, we perform
causal ablations on Qwen2.5-1.5B-Instruct using the 500-sample balanced
SST-2 validation (baseline AUC 0.8488 at this sample size).

\paragraph{Head ablation (12 heads at the final layer).}
We zero each attention head individually at the final transformer layer
by masking the corresponding $128$-dim slice of the o\_proj input.
Across all 12 heads the SST-2 AUC drop is within $\pm 0.002$ of the
baseline (max drop $+0.0019$ for head 3). \emph{No single head carries
the V-axis.}

\paragraph{Block ablation (attention vs MLP).}
We then zero the entire attention block or the entire MLP block at
both paper L27 (the crystallizing layer) and paper L28 (the final
layer):

\begin{table}[h]
\centering
\caption{Block ablation at the last two transformer layers.
SST-2 AUC on 500 balanced samples; baseline is 0.8488. Drops are
all $\leq 0.004$ --- the V-axis is distributed redundantly across
the last two blocks.}
\label{tab:sup_block_ablation}
\begin{tabular}{lcc}
\toprule
Ablation & Paper L27 & Paper L28 \\
\midrule
Zero attention              & 0.8518 & 0.8498 \\
Zero MLP                    & 0.8460 & 0.8457 \\
Zero attention + MLP        & 0.8496 & 0.8480 \\
\bottomrule
\end{tabular}
\end{table}

Even zeroing the entire forward pass of a whole layer (attention + MLP)
changes the final SST-2 AUC by at most $0.004$. Ablating the attention
alone actually \emph{improves} AUC slightly (by 0.003 at L27), possibly
because late-layer attention mixes non-valence features into the residual
stream. Ablating the MLP hurts slightly more, consistent with the MLP
doing the linear re-projection.

\paragraph{Interpretation.}
The final two transformer blocks are \emph{redundant}: either alone
can project the residual stream onto a linear valence direction. This
makes the V-axis a robust convergent feature of late-layer computation
rather than the output of a single ``valence circuit''. It also
explains why the L28 V-axis is nearly orthogonal to the middle-layer
(L14) V-axis --- the late-layer computation builds a direction
$\emph{de novo}$ from distributed features.

\section{Prompt Robustness}
\label{sup:promptrobust}

We tested four different story-generation prompt styles:
(i) \emph{standard} (``Write a short story in which the main character
feels \ldots''), (ii) \emph{first-person} (``Write a first-person diary
entry \ldots''), (iii) \emph{dialogue} (``Write a two-person dialogue in
which one character is clearly feeling \ldots''), and (iv) \emph{scene}
(``Describe a cinematic scene \ldots''). For each style we regenerate 9
stories per class, recompute the V-axis via PCA, and evaluate zero-shot
SST-2 AUC.

\begin{table}[h]
\centering
\caption{Prompt-style robustness on SST-2 (Qwen2.5-1.5B, L28).}
\label{tab:sup_prompt_robust}
\begin{tabular}{lcc}
\toprule
Prompt style & SST-2 AUC & SST-2 acc (median thresh) \\
\midrule
standard      & 0.794 & 0.718 \\
first-person  & 0.761 & 0.677 \\
dialogue      & 0.766 & 0.697 \\
scene         & 0.803 & 0.734 \\
\bottomrule
\end{tabular}
\end{table}

All four styles yield AUC $> 0.75$, confirming the method is robust to
surface-level prompt changes. However, the V-axes are not all mutually
aligned: the pairwise cosine similarities between standard,
first-person, and scene axes are large (0.78--0.82, all positive),
while the \emph{dialogue} axis flips sign relative to the others
($\cos = -0.49$ vs.\ standard, $-0.69$ vs.\ first-person). Dialogue
still yields a strong SST-2 AUC after sign-alignment; the flipped sign
reflects PCA sign ambiguity rather than a genuine semantic disagreement.
This motivates a simple post-hoc \emph{sign-align} step: when deploying a
new prompt style, check the centroid of the \texttt{happy} class lies on
the positive side of the axis, and multiply by $-1$ if not.

\section{Qwen3.5-2B Cross-Scaling}
\label{sup:qwen35}

As a robustness check on model generation, we repeat the main
experiments with Qwen3.5-2B (Instruct, 24 layers), representing a
newer-generation model at a similar parameter scale:

\begin{table}[h]
\centering
\caption{Qwen3.5-2B vs.\ Qwen2.5-1.5B on the main Paper 1 benchmarks.
Qwen2.5-1.5B remains the best single model.}
\label{tab:sup_qwen35}
\begin{tabular}{lcc}
\toprule
Metric & Qwen2.5-1.5B (main) & Qwen3.5-2B \\
\midrule
SST-2 AUC           & \textbf{0.868} & 0.822 \\
SST-2 acc (median)  & \textbf{0.778} & 0.748 \\
IMDB AUC            & \textbf{0.895} & 0.822 \\
EmoBank V $|r|$     & 0.480          & \textbf{0.520} \\
Hu \& Liu $|r|$     & \textbf{0.761} & 0.722 \\
\bottomrule
\end{tabular}
\end{table}

Qwen3.5-2B is competitive but slightly weaker than Qwen2.5-1.5B on 4 of
5 benchmarks. It is \emph{stronger} on the continuous EmoBank regression
(|r| 0.52 vs.\ 0.48). These results confirm that our main findings are
not a Qwen2.5-specific artifact.

\section{Scope Boundaries}
\label{sup:scope_boundaries}

\subsection{Toxicity concept}
We attempted to generalize the method by generating 15 toxic and 15 safe
stories, computing a T-axis as the difference of class means, and testing on
Jigsaw Civil Comments. Balanced AUC was 0.59 (near random). The centroid
distance $|\mu_{\mathrm{toxic}} - \mu_{\mathrm{safe}}| = 58.85$ in L28 space
is large, so the LLM distinguishes its own toxic and safe stories cleanly,
but the direction does not transfer to real-world toxic comments.

\subsection{Multi-LLM ensemble}
Combining V-axes from 4 LLMs and averaging sign-aligned projections:

\begin{table}[h]
\centering
\caption{Multi-LLM ensemble hurts relative to best single LLM.}
\label{tab:sup_multillm}
\begin{tabular}{lccc}
\toprule
Benchmark & Best single (Qwen-1.5B) & Mean 4-model ensemble & $\Delta$ \\
\midrule
SST-2 AUC & 0.837 & 0.805 & $-$0.032 \\
IMDB AUC & 0.819 & 0.779 & $-$0.040 \\
Yelp AUC & 0.928 & 0.907 & $-$0.021 \\
\bottomrule
\end{tabular}
\end{table}

Interpretation: the LLMs have different V-axis qualities (Qwen-1.5B best,
Bloom-560m weakest) and averaging drags the strong model toward the weak
model. Unlike same-model ensembling (which is the core mechanism of our
companion paper on the trajectory-length ensemble trick), cross-model
ensembling suffers from a ``weakest-link'' effect.

\section{Hyperparameters}
\label{sup:hparams}

\paragraph{Story generation.} 9 stories per class $\times$ 9 classes $= 81$
total. Generation: $\max\text{new tokens}=150$, $T=0.8$, top-$p=0.9$,
repetition penalty = 1.1.

\paragraph{Feature extraction.} Layer $L=28$ for Qwen2.5-1.5B.
Mean-pooled over non-pad tokens with attention mask weighting.
Max context = 512 tokens for stories, 128--256 for downstream test text.

\paragraph{PCA.} Standard SVD of the mean-centered 9$\times 1{,}536$
centroid matrix. First principal component taken as V-axis.

\paragraph{Threshold calibration.} For SST-2 calibrated accuracy we use
200 held-out SST-2 samples to find the optimal threshold by grid search.
The ``median'' threshold uses the median of the validation projections
with no labels at all.
